% Self-contained article: text, TikZ figure, and bibliography are included.
% Compile with pdflatex main.tex twice; no BibTeX step is required.
\documentclass[a4paper,11pt]{article}

% Encoding, language, and standard article typography.
\usepackage[T1]{fontenc}
\usepackage[utf8]{inputenc}
\usepackage[english]{babel}
\usepackage{lmodern}
\usepackage[a4paper,margin=25mm,headheight=14pt,headsep=8mm,footskip=12mm]{geometry}
\usepackage{microtype}
\usepackage{amsmath,amssymb,amsthm}
\usepackage{graphicx}
\usepackage{tikz}
\usetikzlibrary{arrows.meta}
\usepackage{caption}
\captionsetup{font=small,labelfont=bf,justification=justified,singlelinecheck=false}
\usepackage{fancyhdr}
\usepackage{xcolor}
\usepackage[numbers,sort&compress]{natbib}
\usepackage{xurl}
\definecolor{linkblue}{RGB}{20,64,136}
\definecolor{figureblue}{RGB}{38,99,145}
\usepackage[colorlinks=true,linkcolor=linkblue,citecolor=linkblue,urlcolor=linkblue]{hyperref}
\newcommand{\doi}[1]{\href{https://doi.org/#1}{doi:\,\nolinkurl{#1}}}
\renewcommand{\bibfont}{\small}
\setlength{\bibsep}{4pt}

% Author information: edit these macros to reuse the template.
\newcommand{\AuthorName}{Osvaldo L. Santos-Pereira}
\newcommand{\AuthorAffiliation}{Physics Institute of Universidade Federal do Rio de Janeiro (UFRJ)}
\newcommand{\AuthorORCID}{0000-0003-2231-517X}
\newcommand{\AuthorEmail}{olsp@if.ufrj.br}
\newcommand{\AuthorFootnote}{%
  ORCID: \href{https://orcid.org/\AuthorORCID}{\AuthorORCID}.\enspace
  E-mail: \href{mailto:\AuthorEmail}{\AuthorEmail}.}

% Mathematical notation.
\newcommand{\R}{\mathbb{R}}
\newcommand{\dd}{\mathop{}\!\mathrm{d}}
\newcommand{\vect}[1]{\boldsymbol{#1}}
\DeclareMathOperator{\tr}{tr}
\numberwithin{equation}{section}
\newtheorem{theorem}{Theorem}[section]

% Page layout: restrained running header and centered page number.
\pagestyle{fancy}
\fancyhf{}
\fancyhead[L]{\small\itshape Liouville's formula}
\fancyhead[R]{\small\itshape Linear differential systems}
\fancyfoot[C]{\thepage}
\renewcommand{\headrulewidth}{0.4pt}
\renewcommand{\footrulewidth}{0.4pt}
\fancypagestyle{plain}{%
  \fancyhf{}
  \fancyfoot[C]{\thepage}
  \renewcommand{\headrulewidth}{0pt}
  \renewcommand{\footrulewidth}{0.4pt}}
\setlength{\parindent}{1.4em}
\setlength{\parskip}{0pt}
\setlength{\emergencystretch}{1.5em}
\clubpenalty=10000
\widowpenalty=10000
\displaywidowpenalty=10000
\raggedbottom

\title{A Geometric Interpretation\\of Liouville's Formula}
\author{\AuthorName\thanks{\AuthorFootnote}\\[0.45em]
  \small\AuthorAffiliation}
\date{}
\hypersetup{pdftitle={A Geometric Interpretation of Liouville's Formula},
  pdfauthor={Osvaldo L. Santos-Pereira},
  pdfsubject={Undergraduate lecture note on linear systems and volume change}}

\begin{document}
\maketitle

\noindent
The solutions of a linear differential system define a family of linear
transformations of the space of initial conditions. This note explains how
the determinant of the fundamental matrix measures the resulting change in
volume and why its rate of change is governed by the trace of the coefficient
matrix. An elementary derivation of Liouville's formula is followed by a
geometric interpretation and two planar examples. The prerequisites are
basic linear algebra, multivariable calculus, and first-order differential
equations.

\section{Introduction}
\label{sec:introduction}

A differential equation describes the motion of an individual point once an
initial condition has been chosen. It also describes the motion of an entire
region when all its points are allowed to evolve according to the same
equation. For a linear system, this motion is determined by a matrix, so the
geometric effect can be studied using determinants. The question addressed
here is how the volume of a region changes as it moves under the system.

The relevant result is usually called \emph{Liouville's formula}, and is also
known as Abel's identity in the theory of linear differential
equations~\cite[Sec.~3.4]{Teschl2012}. It relates the determinant of a
fundamental matrix to the trace of the matrix defining the system. The
derivation below follows the geometric route: first describe the evolving
region, then apply the change-of-variables formula, and finally differentiate
the determinant using Jacobi's formula~\cite[p.~1]{Kahan1998}. All arguments
remain within the setting of real linear systems.

\section{Linear systems and evolving regions}
\label{sec:linear}

We begin by relating the solutions of a linear system to a transformation
of its initial conditions. This requires fixing the initial time and choosing
the normalization of the fundamental matrix explicitly.

Let $I\subseteq\R$ be an interval, let $t_0\in I$, and let $A(t)$ be a real
$n\times n$ matrix whose entries are continuous on $I$. The state
$\vect{x}(t)\in\R^n$ satisfies the initial-value problem
\begin{equation}
  \frac{\dd\vect{x}}{\dd t}=A(t)\vect{x}(t),
  \qquad \vect{x}(t_0)=\vect{x}_0.
  \label{eq:system}
\end{equation}
The existence and uniqueness theorem for linear systems guarantees a unique
solution on $I$ for each initial vector $\vect{x}_0$~\cite[Sec.~3.4]{Teschl2012}.
A fundamental matrix has linearly independent solution vectors as its
columns. Writing $I_n$ for the identity matrix, we choose the normalized
fundamental matrix $U(t)$ defined by
\begin{equation}
  U'(t)=A(t)U(t),
  \qquad U(t_0)=I_n,
  \label{eq:fundamental}
\end{equation}
Its columns are the solutions starting
from the standard basis vectors. Linearity then gives the solution with
arbitrary initial data as
\begin{equation}
  \vect{x}(t)=U(t)\vect{x}_0.
  \label{eq:solution}
\end{equation}
The condition $U(t_0)=I_n$ is what makes $\vect{x}_0$ the actual initial
position in Eq.\,\eqref{eq:solution}. For an unnormalized fundamental matrix
$Y(t)$, the corresponding expression would be
$\vect{x}(t)=Y(t)Y(t_0)^{-1}\vect{x}_0$.

The matrix $U(t)$ is invertible at every time in $I$. Indeed, if
$U(t_*)\vect{c}=\vect{0}$ for some $t_*\in I$, the solution
$U(t)\vect{c}$ passes through the origin at $t_*$. Uniqueness implies that
this is the zero solution throughout $I$, and evaluation at $t_0$ gives
$\vect{c}=\vect{0}$. Thus $U(t)$ has a trivial kernel. Its determinant is
continuous and never vanishes; since $\det U(t_0)=1$, it is positive
throughout $I$.

At a fixed time, the velocity field is $\vect{v}(t,\vect{x})=A(t)\vect{x}$.
Along a solution, its value is $A(t)U(t)\vect{x}_0$, which equals the tangent
vector $\vect{x}'(t)$. If $A$ is constant, this field is independent of time.
The map in Eq.\,\eqref{eq:solution} is a linear deformation of the space of
initial conditions. It fixes the origin and can stretch, compress, rotate,
or shear a region.

Let $\Omega_0\subset\R^n$ be a bounded region with a piecewise smooth boundary
and positive volume. Its image at time $t$ and its volume are denoted by
\begin{equation}
  \Omega_t=U(t)\Omega_0,
  \qquad V(t)=\int_{\Omega_t}\dd^n\vect{x},
  \qquad V_0=V(t_0),
  \label{eq:region}
\end{equation}
where $\dd^n\vect{x}=\dd x_1\cdots\dd x_n$ is the usual volume element.
For $n=2$, the word volume refers to area. The Jacobian matrix of the map
$\vect{x}_0\mapsto U(t)\vect{x}_0$ is $U(t)$ itself. Hence the
change-of-variables formula gives
\begin{equation}
  \begin{aligned}
    V(t)
    &=\int_{\Omega_0}|\det U(t)|\,\dd^n\vect{x}_0\\
    &=\det U(t)\,V_0.
  \end{aligned}
  \label{eq:volume-determinant}
\end{equation}
The absolute value is required in the change-of-variables formula. It can
be removed here because the normalized fundamental matrix has positive
determinant. Since $U(t)$ does not depend on the initial point, its
determinant is the same volume factor everywhere in the region.

\section{Liouville's formula}
\label{sec:liouville}

Equation\,\eqref{eq:volume-determinant} reduces the volume problem to a
scalar quantity. We write $\varphi(t)=\det U(t)$ and derive a differential
equation for $\varphi$. The matrix operation that appears in this equation
is the trace. For an $n\times n$ matrix $B=(b_{ij})$, it is defined by
\begin{equation}
  \tr B=\sum_{i=1}^{n}b_{ii}.
  \label{eq:trace}
\end{equation}
For two $n\times n$ matrices $B$ and $C$, summing their diagonal entries
shows that $\tr(BC)=\sum_{i,j}b_{ij}c_{ji}=\tr(CB)$. This is the cyclic
property of the trace; it does not assert that $BC=CB$.

To differentiate $\varphi$, let $C_{ij}(t)$ denote the cofactor of the entry
$U_{ij}(t)$. Cofactor expansion gives
$\partial(\det U)/\partial U_{ij}=C_{ij}$, while the inverse-matrix formula
gives $C_{ij}=\det U\,(U^{-1})_{ji}$. Applying the chain rule therefore yields
\begin{equation}
  \begin{aligned}
    \varphi'(t)
    &=\sum_{i,j=1}^{n}C_{ij}(t)U'_{ij}(t)\\
    &=\varphi(t)\sum_{i,j=1}^{n}(U(t)^{-1})_{ji}U'_{ij}(t)\\
    &=\varphi(t)\,\tr\!\left(U(t)^{-1}U'(t)\right).
  \end{aligned}
  \label{eq:jacobi}
\end{equation}
This identity is Jacobi's formula~\cite[p.~1]{Kahan1998}. The preceding
invertibility argument justifies the use of $U(t)^{-1}$.

Substituting $U'=AU$ from Eq.\,\eqref{eq:fundamental} makes the differential
equation for the determinant explicit. The cyclic property gives
$\tr(U^{-1}AU)=\tr(AUU^{-1})=\tr A$, so Eq.\,\eqref{eq:jacobi} becomes
\begin{equation}
  \varphi'(t)=\tr A(t)\,\varphi(t),
  \qquad \varphi(t_0)=1.
  \label{eq:scalar-determinant}
\end{equation}
Because $\varphi(t)>0$, we may divide by it and integrate. This gives
\begin{equation}
  \log\varphi(t)-\log\varphi(t_0)
  =\int_{t_0}^{t}\tr A(s)\,\dd s.
  \label{eq:log-determinant}
\end{equation}
Exponentiating and recalling Eq.\,\eqref{eq:volume-determinant} completes
the derivation. We state the result with both its matrix and geometric forms.

\begin{theorem}[Liouville's formula]
\label{thm:liouville}
Let $A:I\to\R^{n\times n}$ be continuous, and let $U$ satisfy
Eq.\,\eqref{eq:fundamental}. Then, for every $t\in I$,
\begin{equation}
  \det U(t)=\exp\!\left(\int_{t_0}^{t}\tr A(s)\,\dd s\right).
  \label{eq:liouville}
\end{equation}
For the evolving region in Eq.\,\eqref{eq:region}, the volume consequently
satisfies
\begin{subequations}\label{eq:volume-law}
\begin{align}
  V'(t)&=\tr A(t)\,V(t),\label{eq:volume-rate}\\
  V(t)&=V_0\exp\!\left(\int_{t_0}^{t}\tr A(s)\,\dd s\right).
  \label{eq:volume-integrated}
\end{align}
\end{subequations}
\end{theorem}

Thus the volume can be determined from the trace without finding every
entry of the fundamental matrix. For a constant matrix $A$, the integral
reduces to $(t-t_0)\tr A$. The same derivation covers continuous
time-dependent matrices.

\section{Geometric meaning of the trace}
\label{sec:geometry}

The trace has a direct interpretation as the instantaneous relative rate
of volume change. This follows by dividing Eq.\,\eqref{eq:volume-rate}
by the positive quantity $V(t)$. A positive trace gives increasing volume
at that time, and a negative trace gives decreasing volume. If
$\tr A(t)=0$ for every $t\in I$, every evolving region preserves its
volume. Conversely, if a region of positive volume has constant volume
throughout $I$, Eq.\,\eqref{eq:volume-rate} implies that the trace vanishes
throughout $I$.

The same quantity appears in elementary vector calculus. If
$A(t)=(a_{ij}(t))$, the components of the velocity field are
$v_i(t,\vect{x})=\sum_{j=1}^{n}a_{ij}(t)x_j$. Its divergence, defined as
the sum of the derivatives of each component with respect to the matching
coordinate, is
\begin{equation}
  \nabla_{\vect{x}}\!\cdot\vect{v}(t,\vect{x})
  =\sum_{i=1}^{n}\frac{\partial v_i}{\partial x_i}
  =\sum_{i=1}^{n}a_{ii}(t)
  =\tr A(t).
  \label{eq:divergence}
\end{equation}
For this linear field, the divergence is spatially uniform, so the same
relative volume change occurs throughout the region. A zero trace at a
single instant only makes the volume derivative vanish at that instant;
volume preservation over an interval requires it to vanish throughout the
interval. Volume preservation allows lengths and angles to change, as the
following examples demonstrate.

\section{Two elementary examples}
\label{sec:examples}

Planar systems make the determinant visible as an area factor. We set
$t_0=0$ and use the unit square $\Omega_0=[0,1]^2$ as the initial region.
The first example allows independent changes along the coordinate axes;
the second gives a shear. Both can be solved by direct integration.

\subsection{Expansion and compression}
\label{sec:diagonal}

Let $a$ and $b$ be real constants and consider the equations $x'=ax$ and
$y'=by$. The coefficient matrix and the normalized fundamental matrix are
\begin{equation}
  A=\begin{pmatrix}a&0\\0&b\end{pmatrix},
  \qquad U(t)=\begin{pmatrix}e^{at}&0\\0&e^{bt}\end{pmatrix}.
  \label{eq:diagonal}
\end{equation}
The unit square becomes a rectangle with side lengths $e^{at}$ and
$e^{bt}$. Its area is therefore
\begin{equation}
  V(t)=e^{(a+b)t},
  \qquad \tr A=a+b.
  \label{eq:diagonal-area}
\end{equation}
This agrees with Eq.\,\eqref{eq:volume-integrated}. The area increases
when $a+b>0$, decreases when $a+b<0$, and remains constant when $a+b=0$.
For the choice $a=1$ and $b=-1$, expansion along the horizontal direction
exactly compensates compression along the vertical direction. At
$t=\log 2$, the sides have lengths $2$ and $1/2$, so the area remains
equal to one. Figure~\ref{fig:deformations}(a) shows this deformation.

\subsection{Shear with constant area}
\label{sec:shear}

Let $k$ be a real constant and consider $x'=ky$ and $y'=0$. Since the
vertical coordinate stays fixed, direct integration gives
$y(t)=y_0$ and $x(t)=x_0+kty_0$. In matrix form,
\begin{equation}
  A=\begin{pmatrix}0&k\\0&0\end{pmatrix},
  \qquad U(t)=\begin{pmatrix}1&kt\\0&1\end{pmatrix}.
  \label{eq:shear}
\end{equation}
Here $\tr A=0$ and $\det U(t)=1$, so the area is constant. The square
becomes a parallelogram whose vertices are $(0,0)$, $(1,0)$,
$(1+kt,1)$, and $(kt,1)$. Its base and height both have length one,
although its sloping sides change length. Figure~\ref{fig:deformations}(b)
shows the case $kt=1$.

\begin{figure}[htbp]
  \centering
  \begin{tikzpicture}[x=1.85cm,y=1.85cm,>=Stealth,font=\small]
    \begin{scope}
      \path[use as bounding box] (-.35,-.55) rectangle (2.65,1.6);
      \fill[figureblue!12] (0,0) rectangle (2,.5);
      \draw[figureblue,line width=.9pt] (0,0) rectangle (2,.5);
      \draw[gray!80,densely dashed,line width=.8pt]
        (0,0) rectangle (1,1);
      \draw[->,black!70] (-.2,0)--(2.4,0) node[right] {$x$};
      \draw[->,black!70] (0,-.18)--(0,1.35) node[above] {$y$};
      \node[below left] at (0,0) {$0$};
      \foreach \x in {1,2}{\draw (\x,.035)--(\x,-.035) node[below] {$\x$};}
      \draw (.035,1)--(-.035,1) node[left] {$1$};
      \draw (.035,.5)--(-.035,.5) node[left] {$\tfrac12$};
      \node[figureblue] at (1.55,.25) {$\Omega_t$};
      \node[gray!80] at (.55,.78) {$\Omega_0$};
      \node[align=center] at (1.1,-.48) {\textbf{(a)} Expansion and compression};
    \end{scope}
    \begin{scope}[xshift=7.35cm]
      \path[use as bounding box] (-.35,-.55) rectangle (2.65,1.6);
      \fill[figureblue!12] (0,0)--(1,0)--(2,1)--(1,1)--cycle;
      \draw[figureblue,line width=.9pt] (0,0)--(1,0)--(2,1)--(1,1)--cycle;
      \draw[gray!80,densely dashed,line width=.8pt]
        (0,0) rectangle (1,1);
      \draw[->,black!70] (-.2,0)--(2.4,0) node[right] {$x$};
      \draw[->,black!70] (0,-.18)--(0,1.35) node[above] {$y$};
      \node[below left] at (0,0) {$0$};
      \foreach \x in {1,2}{\draw (\x,.035)--(\x,-.035) node[below] {$\x$};}
      \draw (.035,1)--(-.035,1) node[left] {$1$};
      \node[figureblue] at (1.4,.65) {$\Omega_t$};
      \node[gray!80] at (.4,.78) {$\Omega_0$};
      \node[align=center] at (1.1,-.48) {\textbf{(b)} Shear};
    \end{scope}
  \end{tikzpicture}
  \caption{Two deformations of the unit square. The dashed outline is
  the initial region $\Omega_0$ and the shaded region is its image
  $\Omega_t$. In (a), $U=\operatorname{diag}(2,1/2)$; in (b),
  $U=\left(\begin{smallmatrix}1&1\\0&1\end{smallmatrix}\right)$.
  Both matrices have determinant one, so both images have unit area.
  The coordinate scales are identical in the two panels.}
  \label{fig:deformations}
\end{figure}

These examples illustrate the two roles of the matrices in the argument.
The fundamental matrix $U(t)$ determines the full deformation, and its
determinant gives the resulting volume factor. The coefficient matrix
$A(t)$ determines the instantaneous velocity field, and its trace gives
the relative rate at which that volume changes.

\begin{thebibliography}{2}

\bibitem[Teschl(2012)]{Teschl2012}
Gerald Teschl.
\newblock \emph{Ordinary Differential Equations and Dynamical Systems}, volume
  140 of \emph{Graduate Studies in Mathematics}.
\newblock American Mathematical Society, Providence, RI, 2012.
\newblock ISBN 978-0-8218-8328-0.
\newblock \doi{10.1090/gsm/140}.
\newblock URL \url{https://www.mat.univie.ac.at/~gerald/ftp/book-ode/ode.pdf}.

\bibitem[Kahan(1998)]{Kahan1998}
William Kahan.
\newblock {Jacobi}'s formula for the derivative of a determinant.
\newblock Lecture notes for Math H110, University of California, Berkeley,
  October 1998.
\newblock URL
  \url{https://people.eecs.berkeley.edu/~wkahan/MathH110/jacobi.pdf}.
\newblock 26 October 1998.

\end{thebibliography}
\end{document}
